\chapter{Diskussion}
\label{ch:Diskussion}

\section{Zusammenfassung und Analyse der Messwerte}
\label{sec:Zusammenfassung}
In diesem Versuch wurde der Stirnling-Motor im Betrib verschiedener Modi untersucht. Insbesondere wurden die Wirkungsgrade der verschiedenen Modi ermittelt.

Für den ersten Versuchsteil wurde der Motor als Kältemaschine betrieben. Über die Kompensation konnte die Kälteleistung $P_\text{K}$ bestimmt werden. Zudem wurde die Verlustwärme pro Zyklus $Q_\text{zy}$ berechnet. Dieser Verlust wird über Reibungskräfte begründet, welche insebsondere bei großen Drehzahlen auftauchen. Mit diesen Werten konnte schließlich der Wirkungsgrad
\res{\eta_\text{K}}{31}{5}{\%}
bestimmt werden.

Im zwieten Versuchsteil wurde der Motor zunächst abermals als Kältemaschine betrieben, während des Betriebes wurde der Modus dann umgeändert, sodass der Motor dann als Wärmepumpe fungiert hat. Dabei wurde die Temperatur des Wassers beobachtet. Die elektrische Heißleistung der Wärmepumpe ist dabei gößer als die Kälteleistung. Ein Vergleich der beiden Kälteleistungen beider Aufgaben zeigt, dass die Kälteleistung $P_\text{K,2}$ im zweiten Aufgabenteil gerde mal $\approx 8 \%$ der Kälteleistung aus dem ersten Aufgabenteil beträgt. Die Standardabweichung beträgt $31.45\sigma$; dies ist nicht signifikant.
Dies ist jedoch der Erwartung entsprechend, es wird nicht derselbe Wert erwartet. Denn die Leistung wird für einen bestimmten Temperaturbereich angegeben. In der ersten Aufgabe fand der Prozess bei Raumtemperatur statt, während das Wasser in der zweiten Aufgabe bei $0^\circ C$ gefriert und die Temperatur bis $\approx - 23^\circ C$ fällt. Somit ist nur bei der zweiten Aufgabe ein (relevanter) Wärmeaustausch mit der Umgebung gegeben, was einen kontinuirlichen Wärmeaustausch zur folge hat. Gegen die Aufgonommene Wäre muss der Motor gegenanarbeiten, was zur Verringerung der Kälteleiszung führt.

Im letzen Versuchsteil wurde der Apperat als Heißluftmotor betrieben. Es wurde erneut die Menge des Wärmeverlustes $Q_\text{ver}$ bestimmt. Dieser wurde zu 
\res{Q_\text{ver}}{6.9}{0.4}{J}
bestimmt. Die Verluste sind wieder insbesondere über Reibungskräfte zu begründen, aber auch durch den Wärmeaustausch. Die Maschine weist in diesem Betrieb verschiedene Temperatur bereiche auf, welche in Wechselwirkung standen. Es wude auch die elektrische Heizleistung bestimmt, welche jedoch keine Wärmeverluste berücksichtigt und damit vermutlich größer bestimmt wurde, als der tatsächliche Wert. 
Es ergibt sich ein Wirkungsgrad von
\res{\eta_\text{u,t}}{6.09}{0.10}{\%}

Dieser Wert liegt deutlich unter dem Wirkungsgrad aus Aufgabe eins.

Zuletzt wurde derselbe Versuchsaufbau noch unter verschiedenen Belastungsgraden betrieben. Es wurden analog zu den vorherigen Aufgaben die Werte berechnet. 
Durch die Last kann zusätzlich ein effektiver Wirkungsgrad bestimmt werden. Die Ergebnisse aller Widerstände sind der \ctab{messdaten} zu entnehmen.
Effektiver und thermischer Wirkungsgrad wurden dann gegen die Frequenz geplottet (\cabb{Wirkungsgrade}). Es lassen sich eindeutge Trends feststellen. Eine Erklärung der Trends ist der \csec{A4} zu entnehmen. 

Es sollte zu letzt noch betont werden, dass es sich bei dem Versuchsaufbau um keinen industriellen Stirlingmotor handelt. Optimierte Stirlingmotoren können dabei höhre Leistungen erbringen. Der Motor im Versuch wurde nicht auf Leistung oder Wirkungsgrad, sondern auf Visualisierung der Funktionsweise optimiert. Dies geht besonders auf die Kosten der Wärmeisolation, denn Glas isoliert nicht gut. 


% \section{Analyse der Messwerte}
% \label{sec:Analyse}

% \section{Kritik}
% \label{sec:Kritik}
